\documentclass[10pt]{article}
\usepackage{setspace}
\onehalfspacing
\usepackage{fontawesome}
\usepackage[newfloat]{minted}
\usepackage{algpseudocode}
\usepackage{algorithm}
\usepackage{mdframed}
\usepackage[margin=1.2in]{geometry}
\usepackage{amsmath,amsthm,amssymb}
\usepackage{mathtools}
\usepackage[colorlinks=true]{hyperref}
\theoremstyle{definition}
\newtheorem{question}{{\normalfont \faGears}~Question}
\newtheorem{compexercise}{{\normalfont \faLaptop}~Implementation exercise}
\theoremstyle{remark}
\newtheorem*{protosolution}{Solution}
\usepackage{xcolor}
\renewcommand{\d}{\mathrm d}


\renewcommand{\mymarks}[1]{\textbf{#1 marks}}
\usepackage{fancyhdr}
\pagestyle{fancy}
\fancyhf{}
\chead{Midterm exam -- Solutions}
\rhead{Page \thepage}

\definecolor{lightgreen}{HTML}{f1faf8}
\newenvironment{solutionframe}
{%
    \begin{mdframed}[
        leftmargin=1cm,
        skipabove=.3cm,
        linecolor=blue,
        backgroundcolor=lightgreen,
        linewidth=0pt,
        innerleftmargin=.5em,
        innerrightmargin=.5em,
        innertopmargin=.3em,
        innerbottommargin=.6em,
    ]
}
{
    \end{mdframed}
}

\newenvironment{solution}
{\pushQED{\qed}\renewcommand{\qedsymbol}{$\triangle$}
\begin{solutionframe}\small \begin{protosolution}}
{\popQED\end{protosolution}\end{solutionframe}}

\begin{document}

\title{Numerical Analysis: Midterm {\small (\mymarks{50})} \\ Solutions}
\author{Urbain Vaes}
\date{Fall 11111101001}
\maketitle

\begin{question}
    [Floating point arithmetic, \mymarks{10}]
    True or false?
    \begin{enumerate}
        \item Let $(\placeholder)_2$ denote binary representation.
            Then
            \(
                (0.1101)_2 + (0.0011)_2 = (1.0)_2.
            \)

        \item
            Let $\varepsilon_{64}$ denote the machine epsilon for the \julia{Float64} format,
            i.e.\ \julia{eps(Float64)}.
            Then the number~2 is representable exactly in this format,
            and the next representable number is $2 + \varepsilon_{64}$.

        \item It holds that
            \(
                (10000)_2 \times (0.1010)_2 = (1010)_2.
            \)
        \item
            In Julia, \julia{Float64(.6) == Float32(.6)} evaluates to \julia{true}.

        \item
            The spacing (in absolute value) between successive single-precision (\julia{Float32}) floating point numbers is constant.

        \item
            Infinitely many distinct real numbers can be represented exactly in the~\julia{Float64} format,
            but only finitely many can be represented exactly in the~\julia{Float32} format.

        \item
            It holds that $(0.\overline{101})_2 = \frac{5}{7}$.

        \item
            Machine addition~$\madd$ is an operation that is \emph{associative} but not \emph{commutative}.

        \item
            The machine epsilon is the smallest number of the form~$2^{-n}$ with $n \in \nat$ that can be represented exactly in a floating point format.

        \item
            In Julia, the expression \julia{1 + eps()/3 == 1 + eps()} evaluates to \julia{true}.

        \item
            \textbf{Bonus.}
            In Julia, the expression \julia{exp(eps()/2) == 1 + eps()} evaluates to \julia{true}.

        \item
            \textbf{Bonus.}
            In Julia, the expression \julia{cos(eps()) == 1} evaluates to \julia{true}.
    \end{enumerate}
\end{question}

\begin{solution}
    Below, $\varepsilon = \varepsilon_{64} = 2^{-52}$.
    \begin{enumerate}
        \item \textbf{True}:
            $(0.1101)_2 + (0.0011)_2 = \frac{13}{16} + \frac{3}{16} = 1$.

        \item \textbf{False}:
            the number 2 is representable as $2^1 \times (1.00\dots0)_2$,
            but the spacing between successive \julia{Float64} numbers in $[2, 4)$ is $2^1 \times 2^{-52} = 2 \varepsilon$,
            so the next representable number is $2 + 2\varepsilon$.

        \item \textbf{True}:
            multiplying by $(10000)_2 = 2^4$ shifts the binary point by four positions to the right,
            so $(10000)_2 \times (0.1010)_2 = (1010)_2$.
            In decimal notation, $16 \times 0.625 = 10$.

        \item \textbf{False}:
            the number $0.6 = (0.1\overline{0011})_2$ has an infinite binary expansion,
            so it is rounded to 53 significant bits by \julia{Float64} and to 24 bits by \julia{Float32}.
            The two results are different real numbers,
            and the comparison, which converts the \julia{Float32} number to \julia{Float64} exactly, returns \julia{false}.

        \item \textbf{False}:
            the spacing is $2^{E} \times 2^{-23}$ between numbers with exponent~$E$,
            so it doubles at every power of 2.

        \item \textbf{False}:
            both formats contain only finitely many numbers,
            since they are encoded with finitely many bits (64 and 32, respectively).

        \item \textbf{True}:
            if $x = (0.\overline{101})_2$, then $2^3 x = (101.\overline{101})_2 = 5 + x$,
            so $7x = 5$.

        \item \textbf{False}:
            it is the other way around.
            Machine addition is commutative,
            since $x \madd y = \operatorname{fl}(x + y) = \operatorname{fl}(y + x) = y \madd x$,
            but it is not associative, because intermediate results are rounded.
            For example, in \julia{Float64},
            $(1 \madd \varepsilon/2) \madd \varepsilon/2 = 1$ whereas $1 \madd (\varepsilon/2 \madd \varepsilon/2) = 1 + \varepsilon$.

        \item \textbf{False}:
            much smaller powers of 2 can be represented exactly,
            for example $2^{-1022}$ in \julia{Float64}.
            The machine epsilon is the distance between 1 and the next representable number.

        \item \textbf{False}:
            the exact sum $1 + \varepsilon/3$ is closer to 1 than to $1 + \varepsilon$,
            so \julia{1 + eps()/3} evaluates to~1,
            which is different from \julia{1 + eps()}.

        \item \textbf{True}:
            by Taylor's theorem,
            $\exp(\varepsilon/2) = 1 + \varepsilon/2 + \varepsilon^2/8 + \dotsb$,
            which is slightly larger than $1 + \varepsilon/2$,
            the midpoint between the successive floating point numbers $1$ and $1 + \varepsilon$.
            Rounding to the nearest floating point number therefore gives $1 + \varepsilon$.

        \item \textbf{True}:
            $\cos(\varepsilon) = 1 - \varepsilon^2/2 + \dotsb$.
            The floating point number preceding~1 is $1 - \varepsilon/2$,
            and since $\varepsilon^2/2$ is much smaller than $\varepsilon/4$,
            the nearest floating point number to $\cos(\varepsilon)$ is~1.
    \end{enumerate}
\end{solution}

\begin{question}
    [Interpolation and approximation, \mymarks{10}]
    Are the following statements true or false? Prove or disprove.
    Recall that~$\mathcal P_d$ denotes the set of polynomials of degree at most~$d$.

    \begin{enumerate}
        \item
            Assume that $x_0 < x_1 < x_2 < x_3$ and $y_0, y_1, y_2, y_3$ are given real numbers.
            Then there exists a polynomial~$p \in \mathcal P_3$ such that
            \(
                p(x_i) = y_i
            \)
            for all $i \in \{0, 1, 2, 3\}$.

        \item
            Assume that $x_0 < x_1 < x_2$ and $y_0, y_1, y_2$ are given real numbers.
            Then, there can exist \emph{at most one} polynomial~$p \in \mathcal P_2$ such that
            \(
                p(x_i) = y_i
            \)
            for all $i \in \{0, 1, 2\}$.

        \item
            Let $p \in \mathcal P_d$ be a polynomial of degree $d > 0$,
            and let $q\colon \real \to \real$ be given by $q(x) = p(x+1) - p(x)$.
            Then it holds that~$q \in \mathcal P_{d-1}$.

        \item
            For $n \in \nat$, let $x^n_i = i/n$ for $i = 0, 1, \dotsc, n$.
            Assume that $u \colon \real \to \real$ is the smooth function given by~$u(x) = \sin(3x) + x^3$,
            and let $p_n \in \mathcal P_n$ denote the polynomial interpolation of~$u$ at the points $x^n_0, x^n_1, \dotsc, x^n_n$.
            Then it holds that
            \[
                \lim_{n \to \infty} \left( \max_{x \in [0, 1]} \bigl\lvert u(x) - p_n(x) \bigr\rvert \right) = 0 \, .
            \]
    \end{enumerate}
\end{question}

\begin{solution}
    \begin{enumerate}
        \item \textbf{True}:
            the Lagrange interpolating polynomial
            \[
                p(x) = \sum_{i=0}^{3} y_i \, \ell_i(x),
                \qquad
                \ell_i(x) = \prod_{j \neq i} \frac{x - x_j}{x_i - x_j},
            \]
            belongs to $\mathcal P_3$ and satisfies $p(x_i) = y_i$,
            because $\ell_i(x_j) = 1$ if $i = j$ and 0 otherwise.

        \item \textbf{True}:
            if $p, \widetilde p \in \mathcal P_2$ both satisfy the interpolation conditions,
            then $p - \widetilde p \in \mathcal P_2$ has the three distinct roots $x_0, x_1, x_2$.
            A nonzero polynomial of degree at most 2 has at most 2 roots,
            so $p - \widetilde p = 0$.

        \item \textbf{True}:
            write $p(x) = \sum_{k=0}^{d} a_k x^k$.
            By the binomial formula,
            $(x+1)^k - x^k = \sum_{j=0}^{k-1} \binom{k}{j} x^j$ is a polynomial of degree $k - 1$ for $k \geq 1$,
            and it is zero for $k = 0$.
            Therefore $q(x) = \sum_{k=1}^{d} a_k \bigl((x+1)^k - x^k\bigr)$ belongs to $\mathcal P_{d-1}$.

        \item \textbf{True}:
            by the interpolation error formula,
            for every $x \in [0, 1]$ there exists $\xi \in [0, 1]$ such that
            \[
                u(x) - p_n(x) = \frac{u^{(n+1)}(\xi)}{(n+1)!} \prod_{i=0}^{n} (x - x^n_i).
            \]
            Each factor of the product satisfies $\lvert x - x^n_i \rvert \leq 1$,
            and, for $n \geq 3$, the derivatives of $x^3$ of order $n+1$ vanish,
            so $\lvert u^{(n+1)}(\xi) \rvert = 3^{n+1} \lvert \sin^{(n+1)}(3 \xi) \rvert \leq 3^{n+1}$.
            Consequently,
            \[
                \max_{x \in [0, 1]} \bigl\lvert u(x) - p_n(x) \bigr\rvert
                \leq \frac{3^{n+1}}{(n+1)!}
                \xrightarrow[n \to \infty]{} 0.
            \]
            Note that Runge's phenomenon does not occur here,
            because all the derivatives of~$u$ are bounded by $3^{n+1}$, up to a constant.
    \end{enumerate}
\end{solution}

\begin{question}
    [Interpolation, open-ended question, \mymarks{5}]
    In polynomial interpolation, the error depends both on the function being interpolated and the choice of interpolation nodes. Consider two families of nodes on the interval $[-1,1]$:
    \begin{itemize}
        \item Equally spaced nodes,
        \item Chebyshev nodes.
    \end{itemize}
    Discuss qualitatively (and illustrate with examples or plots if you wish) how the interpolation error behaves as the degree
    $n$ increases in each case.
    Why does one choice of nodes perform better for large~$n$,
    and what is the mathematical motivation for using Chebyshev nodes?
    \emph{You may refer to Runge's phenomenon, but go beyond merely stating it.}
\end{question}

\begin{solution}
    Many answers are possible; a good answer should contain the following elements.
    For a smooth function $u$ and interpolation nodes $x_0, \dotsc, x_n$ in $[-1, 1]$,
    the interpolation error satisfies
    \[
        u(x) - p_n(x) = \frac{u^{(n+1)}(\xi)}{(n+1)!} \, \omega_{n+1}(x),
        \qquad
        \omega_{n+1}(x) = \prod_{i=0}^{n} (x - x_i).
    \]
    The first factor depends only on the function~$u$,
    while the second depends only on the nodes.
    \begin{itemize}
        \item
            \emph{Equally spaced nodes.}
            The function $\omega_{n+1}$ is small in the middle of the interval but very large near the endpoints,
            where its maximum is larger than in the middle by a factor growing exponentially with~$n$.
            If the derivatives of $u$ grow quickly,
            as for the Runge function $u(x) = (1 + 25 x^2)^{-1}$,
            whose $(n+1)$-th derivative grows roughly like $(n+1)! \, 5^{n+1}$,
            the error grows with $n$ near the endpoints:
            the interpolant oscillates with increasing amplitude there (Runge's phenomenon).
            Equivalently, the Lebesgue constant of equally spaced nodes grows like $2^n$,
            so the interpolation can be much worse than the best polynomial approximation,
            and it is also very sensitive to rounding errors and noise in the data.

        \item
            \emph{Chebyshev nodes}
            $x_i = \cos\bigl(\frac{(2i+1)\pi}{2(n+1)}\bigr)$, $i = 0, \dotsc, n$,
            are the roots of the Chebyshev polynomial $T_{n+1}$.
            They cluster near the endpoints,
            precisely where equally spaced nodes perform poorly.
            For these nodes, $\omega_{n+1} = 2^{-n} T_{n+1}$,
            so that $\max_{[-1, 1]} \lvert \omega_{n+1} \rvert = 2^{-n}$,
            and this is the smallest possible value over all choices of $n+1$ nodes:
            Chebyshev nodes minimize the node-dependent factor in the error formula.
            Their Lebesgue constant grows only logarithmically with $n$,
            so the interpolant is close to the best polynomial approximation.
            For the Runge function, or any function analytic in a neighborhood of $[-1, 1]$,
            the error decreases exponentially with $n$.
    \end{itemize}
    A plot of the error $\max_{[-1, 1]} \lvert u - p_n \rvert$ against $n$ for the Runge function
    shows exponential growth for equally spaced nodes and exponential decay for Chebyshev nodes.
\end{solution}

\begin{compexercise}
    [Interpolation, \mymarks{5}]
    Write Julia code that computes and plots the interpolating polynomial~$p \in \mathcal P_3$ through the following points:
    $(0, 0)$, $(1, 4)$, $(2, 15)$, $(3, 40)$.
    The plot should display both the interpolation points and the graph of the interpolating polynomial over an appropriate range.
    Do not use any other library than the ones already imported.
\end{compexercise}

\begin{solution}
    We solve the Vandermonde system for the coefficients of $p(x) = \alpha_0 + \alpha_1 x + \alpha_2 x^2 + \alpha_3 x^3$.
    \begin{minted}{julia}
using LinearAlgebra
using Plots
x = [0, 1, 2, 3]
y = [0, 4, 15, 40]
V = [xi^j for xi in x, j in 0:3]  # Vandermonde matrix
α = V \ y
p(t) = sum(α[j+1] * t^j for j in 0:3)
t = range(-0.5, 3.5, length=200)
plot(t, p.(t), label="Interpolating polynomial")
scatter!(x, y, label="Interpolation points")
    \end{minted}
    The interpolating polynomial is $p(x) = \frac{17}{6} x + \frac{7}{6} x^3$.
\end{solution}

\begin{question}
    [Numerical integration, \mymarks{5}]
    Are the following statements true or false?
    Justify briefly.
    \begin{enumerate}
        \item
            The degree of precision of the following quadrature rule is 2:
            \[
                \int_{-1}^{1} u(x) \, \d x \approx 2u(0) \, .
            \]

        \item
            The degree of precision of the following rule is equal to 3:
            \[
                \int_{-1}^{1} u(x) \, \d x \approx u \left(-\frac{1}{3}\right) + u \left(\frac{1}{3}\right).
            \]

        \item
            For any natural number $N > 0$,
            there exists a quadrature rule with a degree of precision equal to $2N - 1$ of the form
            \[
                \int_{-1}^{1} u(x) \, \d x \approx
                \sum_{n=1}^{N} w_n u(x_n).
            \]

        \item
            Let $x^N_i = i/N$ and consider the following approximation of $\int_{0}^{1} u(x) \, \d x$:
            \begin{equation}
                \label{eq:trapezium}
                \widehat I_N = \frac{1}{2N} \Bigl( u\left(x^N_0\right) + 2 u\left(x^N_1\right) + 2 u\left(x^N_2\right) + \dotsc + 2 u\left(x^N_{N-2}\right) + 2 u\left(x^N_{N-1}\right) + u\left(x^N_N\right) \Bigr).
            \end{equation}
            Suppose first that $u$ is the Runge function, given by $u(x) = (1 + 25 x^2)^{-1}$.
            Then $\widehat I_N$ diverges in the limit $N \to \infty$.

        \item
            Let $u(x) = \cos(3x)$ and let $\widehat I_N$ be as in~\eqref{eq:trapezium}.
            Then it holds that
            \[
                \lim_{N \to +\infty}
                \left( \left\lvert \widehat I_N - \int_{0}^1 u(x) \, \d x \right\rvert \right) = 0.
            \]

        \item
            (\textbf{Bonus.})
            Fix $u(x) = 2x - 1$ and let $\widehat I_N$ be as in~\eqref{eq:trapezium}.
            Then $\widehat I_N = 0$ for all $N \geq 2$.
    \end{enumerate}
\end{question}

\begin{solution}
    \begin{enumerate}
        \item \textbf{False}:
            the rule is exact for $u(x) = 1$ and $u(x) = x$,
            but not for $u(x) = x^2$,
            since $\int_{-1}^{1} x^2 \, \d x = \frac{2}{3}$ whereas $2 \times 0^2 = 0$.
            Its degree of precision is 1.

        \item \textbf{False}:
            the rule is exact for $u(x) = 1$ (both sides equal 2) and $u(x) = x$ (both sides equal 0),
            but for $u(x) = x^2$ the rule gives $\frac{2}{9} \neq \frac{2}{3}$.
            Its degree of precision is 1.
            The rule with degree of precision 3 is the Gauss--Legendre rule with nodes $\pm 1/\sqrt{3}$.

        \item \textbf{True}:
            the Gauss--Legendre rule with $N$ nodes,
            whose nodes are the roots of the Legendre polynomial of degree $N$,
            has degree of precision $2N - 1$.

        \item \textbf{False}:
            $\widehat I_N$ is the composite trapezium rule.
            Since the Runge function is smooth on $[0, 1]$,
            the error satisfies $\bigl\lvert \widehat I_N - \int_0^1 u(x) \, \d x \bigr\rvert \leq \frac{1}{12 N^2} \max_{[0, 1]} \lvert u'' \rvert$,
            so $\widehat I_N$ converges to the integral.
            (Runge's phenomenon concerns interpolation by a single polynomial of high degree,
            not composite rules.)

        \item \textbf{True},
            by the same error bound: the error is at most $\frac{9}{12 N^2}$, since $\lvert u'' \rvert \leq 9$.

        \item \textbf{True}:
            the trapezium rule is exact for polynomials of degree~1 on each subinterval,
            so $\widehat I_N = \int_0^1 (2x - 1) \, \d x = 0$ for all $N \geq 1$.
            This can also be checked directly,
            since $u(x^N_i) + u(x^N_{N-i}) = 0$ and the weights are symmetric.
    \end{enumerate}
\end{solution}

\begin{question}
    [Gaussian--Hermite numerical integration, \mymarks{10}]

The Gauss--Hermite quadrature formula with \(n\) nodes is an approximation of the form
\[
I(u) \coloneq \int_{-\infty}^{\infty} u(x)\, \e^{-x^{2}} \, \d x \approx \sum_{i=1}^n w_i\,u(x_i)=:\widehat I_n(u),
\]
which is exact when \(u\) is a polynomial of degree \(\le 2n-1\).
Note that the nodes are numbered \(1,\dots,n\).
For this question, we take for granted that, for integers \(i\ge 0\),
it holds that
\[
    \int_{-\infty}^{\infty} x^{i} \e^{-x^{2}}\, \d x =
    \begin{cases}
        \displaystyle
        0, & \text{if $i$ is odd},\\[2mm]
        \displaystyle
        (i-1)!! \, \sqrt{\frac{\pi}{2^i}}, & \text{if $i$ is even},
    \end{cases}
\]
where
\(
    (i-1)!! \coloneq 1 \times 3 \times 5 \times \cdots \times (i-1) \, .
\)
In particular, with all the integrals being over $(-\infty, \infty)$,
the following special cases may be useful in your computations:
\[
    \int \e^{-x^{2}}\, \d x = \sqrt{\pi}, \quad
    \int x^2 \, \e^{-x^{2}}\, \d x = \frac{\sqrt{\pi}}{2}, \quad
    \int x^4 \, \e^{-x^{2}}\, \d x = \frac{3}{4} \sqrt{\pi}, \quad
    \int x^6 \, \e^{-x^{2}}\, \d x = \frac{15}{8} \sqrt{\pi}\, .
\]

\begin{enumerate}
\item (\mymarks{5})
Find the nodes and weights of the Gauss--Hermite rule with \(n=3\) nodes.
By symmetry,
we expect nodes of the form $(-z,0,z)$ and weights $(w_1,w_2,w_1)$,
which reduces the number of unknowns to three.

\begin{solution}
    By symmetry, the rule is exact for all odd monomials.
    Requiring exactness for $1$, $x^2$ and $x^4$ gives
    \[
        2 w_1 + w_2 = \sqrt{\pi},
        \qquad
        2 w_1 z^2 = \frac{\sqrt{\pi}}{2},
        \qquad
        2 w_1 z^4 = \frac{3}{4} \sqrt{\pi}.
    \]
    Dividing the third equation by the second gives $z^2 = \frac{3}{2}$,
    and then $w_1 = \frac{\sqrt{\pi}}{4 z^2} = \frac{\sqrt{\pi}}{6}$ and $w_2 = \sqrt{\pi} - 2 w_1 = \frac{2\sqrt{\pi}}{3}$.
    The nodes and weights are therefore
    \[
        (x_1, x_2, x_3) = \left(-\sqrt{\tfrac{3}{2}}, 0, \sqrt{\tfrac{3}{2}}\right),
        \qquad
        (w_1, w_2, w_3) = \left(\frac{\sqrt{\pi}}{6}, \frac{2\sqrt{\pi}}{3}, \frac{\sqrt{\pi}}{6}\right).
    \]
    The rule is exact for all polynomials of degree at most $5 = 2n - 1$, but not for $x^6$:
    it gives $2 w_1 z^6 = \frac{9}{8} \sqrt{\pi} \neq \frac{15}{8} \sqrt{\pi}$.
\end{solution}

\item (\mymarks{5})
    Let $\{H_0, H_1, \dotsc\}$ denote orthogonal polynomials for the inner product
    \[
        \ip{f, g} := \int_{-\infty}^{\infty} f(x) g(x) \, \e^{-x^2} \, \d x
    \]
    which, in addition,
    satisfy the following two conditions:
    \begin{itemize}
        \item
            For all $i \in \nat$, the polynomial $H_i$ is of degree $i$.
        \item
            The leading coefficient of $H_i$, which multiplies $x^i$, is equal to 1.
    \end{itemize}
    Calculate $H_0$, $H_1$, $H_2$ and $H_3$.
    What is the relationship between \(H_3\) and the quadrature rule found in the first item?

\begin{solution}
    Since the weight $\e^{-x^2}$ is even,
    $H_i$ is even when $i$ is even and odd when $i$ is odd,
    so all odd-degree (resp.\ even-degree) inner products vanish automatically.
    We obtain
    \begin{itemize}
        \item $H_0(x) = 1$;
        \item $H_1(x) = x$, since $\ip{x, 1} = 0$;
        \item $H_2(x) = x^2 - a$ with $\ip{x^2 - a, 1} = \frac{\sqrt{\pi}}{2} - a \sqrt{\pi} = 0$, so $H_2(x) = x^2 - \frac{1}{2}$;
        \item $H_3(x) = x^3 - b x$ with $\ip{x^3 - bx, x} = \frac{3}{4}\sqrt{\pi} - b \frac{\sqrt{\pi}}{2} = 0$, so $H_3(x) = x^3 - \frac{3}{2} x$.
    \end{itemize}
    The roots of $H_3(x) = x \bigl(x^2 - \frac{3}{2}\bigr)$ are $0$ and $\pm \sqrt{3/2}$:
    they are exactly the nodes of the Gauss--Hermite rule with 3 nodes.
    This is a general fact: the nodes of the Gauss quadrature rule with $n$ nodes
    are the roots of the orthogonal polynomial of degree $n$ for the corresponding weight.
\end{solution}

\item (Bonus, \mymarks{2})
    Calculate $H_4$ and, using this result,
    deduce the nodes and weights of the Gauss--Hermite quadrature with~$4$ points.

\begin{solution}
    By parity, $H_4(x) = x^4 + b x^2 + c$, and orthogonality to $1$ and $x^2$ gives
    \[
        \frac{3}{4} + \frac{b}{2} + c = 0,
        \qquad
        \frac{15}{8} + \frac{3b}{4} + \frac{c}{2} = 0.
    \]
    Subtracting the first equation from twice the second gives $3 + b = 0$,
    so $b = -3$, $c = \frac{3}{4}$ and
    \(
        H_4(x) = x^4 - 3 x^2 + \frac{3}{4}.
    \)
    Its roots satisfy $x^2 = \frac{3 \pm \sqrt{6}}{2}$,
    so the nodes of the rule are $\pm a$ and $\pm b$ with
    \[
        a = \sqrt{\frac{3 - \sqrt{6}}{2}} \approx 0.5246,
        \qquad
        b = \sqrt{\frac{3 + \sqrt{6}}{2}} \approx 1.6507.
    \]
    By symmetry, the weights are of the form $(w_b, w_a, w_a, w_b)$,
    and the rule is exact for odd polynomials.
    Exactness for $1$ and $x^2$ gives
    $w_a + w_b = \frac{\sqrt{\pi}}{2}$ and $w_a a^2 + w_b b^2 = \frac{\sqrt{\pi}}{4}$,
    whose solution is
    \[
        w_a = \frac{3 + \sqrt{6}}{12} \sqrt{\pi} \approx 0.8049,
        \qquad
        w_b = \frac{3 - \sqrt{6}}{12} \sqrt{\pi} \approx 0.0813.
    \]
    Since the nodes are the roots of $H_4$,
    the rule is then exact for all polynomials of degree at most 7.
\end{solution}
\end{enumerate}
\end{question}

\begin{compexercise}
    [Numerical integration, \mymarks{10}]
    The midpoint quadrature rule reads
    \[
        \int_{-1}^{1} u(x) \, \d x
        \approx 2 u(0) \, .
    \]
    \begin{itemize}
        \item
            (\mymarks{3})
            Write a function \julia{midpoint(u, a, b)} that returns, using this quadrature rule,
            an approximation of the integral
            \begin{align}
                \label{eq:integral_midpoint}
                \int_{a}^{b} u(x) \, \d x \, .
            \end{align}

        \item
            (\mymarks{4})
            Write a function \julia{composite_midpoint(u, a, b, N)}
            that returns an approximation of the integral~\eqref{eq:integral_midpoint},
            this time using a composite version of the midpoint rule.
            More precisely, the approximation should be obtained
            by partitioning the integration interval $[a, b]$ into $N$ cells,
            and applying the midpoint rule within each cell.

        \item
            (\mymarks{3})
            Take $u(x) = \cos(x)$, $a = -1$ and $b = 1$.
            In this setting, plot the evolution of the error for $N$ varying from 1 to 1000.
    \end{itemize}
\end{compexercise}

\begin{solution}
    By the change of variable $x = \frac{a+b}{2} + \frac{b-a}{2} t$,
    the midpoint rule on $[a, b]$ reads $\int_a^b u(x) \, \d x \approx (b - a) \, u\bigl(\frac{a+b}{2}\bigr)$.
    \begin{minted}{julia}
function midpoint(u, a, b)
    return (b - a) * u((a + b) / 2)
end

function composite_midpoint(u, a, b, N)
    h = (b - a) / N
    return sum(midpoint(u, a + (i-1)*h, a + i*h) for i in 1:N)
end

using Plots
u(x) = cos(x)
I = 2sin(1)  # Exact value of the integral
Ns = 1:1000
errors = [abs(composite_midpoint(u, -1, 1, N) - I) for N in Ns]
plot(Ns, errors, xscale=:log10, yscale=:log10,
     xlabel="N", ylabel="Error", label="Composite midpoint rule")
plot!(Ns, Ns.^(-2), linestyle=:dash, label="1/N²")
    \end{minted}
    In logarithmic scale, the error decreases along a line of slope $-2$:
    the composite midpoint rule converges at rate $\mathcal O(N^{-2})$.
    More precisely, the error is approximately $\frac{h^2}{24} \bigl\lvert u'(b) - u'(a) \bigr\rvert = \frac{\sin(1)}{3N^2}$ with $h = 2/N$.
    Note that the code template in the exam statement omits the argument~\julia{N} of \julia{composite_midpoint}.
\end{solution}

\end{document}
